% Section fragment. Requires hyperref; no bibliography keys are used.
\section{Recent literature and research positioning}
\label{sec:literature}
\noindent Search cutoff: 16 September 2026. This focused scoping search covers the
original ASDPose release and nine complementary primary studies. The full search
log, source register, methodological qualifications, and download manifest are
preserved in \texttt{reports/literature\_survey.md} and
\texttt{references/papers/literature\_downloads.json}.

\subsection{Has the released dataset been reused?}
No independent reuse of the released ASDPose files was verified in this search.
This is a bounded retrieval result, not evidence that no reuse exists. Exact
ASDMotion/ASDPose searches and citation-neighborhood searches found relevant
studies with different cohorts. For example, Freud et al. and Amraee et al. cite
Barami but describe separate experiments. Shared authors or an ADOS assessment
context do not establish identical data. The original release remains the direct
baseline: \href{https://doi.org/10.1001/jamanetworkopen.2024.32851}{Barami et al.,
JAMA Network Open, 2024}.

\subsection{Different questions require different observations}
\begin{description}
\item[Repetition structure.]
\href{https://doi.org/10.1038/s41598-026-60095-8}{MBaye et al., Scientific Reports,
July 2026} construct interpretable persistent-homology features. The publisher
identifies an accepted early article; detailed methods in our review refer to
the \href{https://doi.org/10.1101/2025.09.03.674008}{September 2025 preprint}.
It uses a distinct six-participant video/accelerometer cohort. Periodicity is
a useful candidate representation, not demonstrated ASDPose transfer.

\item[Concurrent movement categories.]
\href{https://doi.org/10.1002/aur.70020}{Lemler et al., Autism Research, March 2025}
use OpenPose and an LSTM for multi-label flapping/jumping recognition. With
subject-disjoint nested validation, macro accuracy is 0.702 but macro F1 is
0.318. This study makes performance under new-child evaluation especially
important; its Frankfurt cohort is separate.

\item[Multimodal action recognition.]
\href{https://doi.org/10.1109/JBHI.2024.3511601}{Zhang et al., IEEE JBHI,
March 2025 issue} combine RGB, optical flow, and skeletal structure with APMFNet
on ACSA653. Such fusion requires modalities that cannot be recovered from a
body-skeleton release. The author manuscript is archived locally for protocol
comparison.

\item[Interpersonal coordination.]
\href{https://doi.org/10.1038/s41598-024-56098-y}{Koehler et al., Scientific
Reports, March 2024} classify patient--clinician dyads from motion synchrony.
Balanced accuracy is 63.4\% against clinical controls. The authors discuss the
possible influence of clinician movement. This is a dyadic clinical-comparison
question rather than SMM episode detection.

\item[Neural mechanisms of imitation.]
\href{https://doi.org/10.1038/s41597-025-05313-0}{Paulo et al., Scientific Data,
June 2025} introduce Move4AS: EEG with marker-based 3D motion during walking
and dancing imitation. It enables neural--motor questions requiring different
sensors and an experimental task.

\item[Purposeful fine motor control.]
\href{https://doi.org/10.1002/aur.70049}{Freud et al., Autism Research, May 2025}
analyze thumb/index-finger grasping in young adults, obtaining over 84\%
classification accuracy with subject-wise validation. The population and
fine-grained measurement differ from preschool body-pose video. Their result
does not validate diagnosis from ASDPose.

\item[Measurement agreement and facial context.]
\href{https://doi.org/10.1186/s13229-025-00685-x}{Manelis-Baram et al., Molecular
Autism, October 2025} compare three facial-expression tools. Outputs differ,
but expression quantity does not significantly distinguish groups. The study
leaves timing, quality, and social context unresolved and is not verified reuse
of the released skeleton data.

\item[Event-level kinematics.]
\href{https://doi.org/10.1111/desc.70176}{Stenum et al., Developmental Science,
March 2026 online} find that event-level amplitude differences disappear after
averaging. This separate toddler cohort cites Barami but does not reuse ASDPose.
Data access requires author request and a data-use agreement.

\item[Semantic behavior context.]
\href{https://openaccess.thecvf.com/content/CVPR2026W/CV4Smalls/papers/Amraee_Toward_Automated_Behavior_Understanding_in_Autism_A_Zero-Shot_Vision-Language_Model_CVPRW_2026_paper.pdf}{Amraee
et al., CVPR Workshops, June 2026} classify 40 clips with vision-language
descriptions and prompted language models. The target includes self-injury,
aggression, and property destruction. This requires visual context and is a
small separate evaluation, not ASDPose benchmarking.
\end{description}

\subsection{Proposed research direction: trustworthy motion measurement}
The following proposals are our synthesis, not results reported by the cited
studies. A useful first research question is:
\begin{quote}
How reliably can SMM burden be estimated for unseen children when pose quality,
missing intervals, and temporal uncertainty are explicitly represented?
\end{quote}
This question has a data-availability gate. The inspected training release
contains 36,930 clips; it does not establish a complete recording timeline or
continuous annotation coverage. Boundary units, source-recording correspondence,
and participant mapping require confirmation. Consequently, clip-level
measurement reliability is the directly feasible first stage. Full-session
burden evaluation must wait for validated continuous recordings, complete event
intervals, and an observed-time denominator. These are local audit findings,
documented in \texttt{data/processed/dataset\_summary.json},
\texttt{dataset\_structure.json}, and \texttt{eda\_notes.json}; the continuous
archive was pending at this review stage. The proposed sequence is:
\begin{enumerate}
\item Freeze participant-disjoint splits before choosing features or thresholds.
Verify interval units and continuous label coverage before estimating burden.
Summarize recording length, repeated sessions, visible-body coverage, joint
confidence, missing runs, and track discontinuities.
\item Establish simple motion-energy and periodicity baselines. Compare
autocorrelation/spectral features with recurrence or topological features before
claiming benefits from a more complex representation.
\item Reproduce the original temporal labeling and inference rules. Report all
deviations between paper, repository, and available release.
\item Compare confidence-aware features or masking and multiple temporal scales
with the same baseline, children, labels, and evaluation procedure.
\item Evaluate precision--recall, episode matching, false events per hour, and
duration/burden error. Bootstrap children, not overlapping windows. Tune
calibration and thresholds on validation data only.
\item Explore clinical associations only after reliability is demonstrated.
Handle repeated sessions and prespecify covariates; do not interpret association
as a causal mechanism or treatment response.
\end{enumerate}

\paragraph{Tests that could disprove the proposed benefit.}
An uncertainty-aware method is not useful merely because its aggregate score
improves. Check whether the improvement persists at matched recall, in children
with fewer visible joints, and after matching movement energy. Test ordinary
periodic activity, track jumps, weak/aperiodic events, overlapping people, and
short bouts as separate error strata. Report abstention coverage and the share
of time that cannot be measured. Negative results are informative if they expose
a representation limit.

\paragraph{Proposed event-distribution analysis.}
Motivated by Stenum's distinction between events and averages, compare
within-child distributions of wrist/body amplitude and recurrence instead of
only a single child mean. This is our proposed experiment, not a published
ASDPose result. The release includes subtype and combined action strings, which
need an ontology and concurrency audit. A label such as ``Fingers'' does not
provide finger coordinates in the observed 17-joint body representation.

\paragraph{Questions requiring additional data.}
Diagnostic discrimination needs an appropriate non-ASD clinical comparison
cohort. Subtype discovery needs trustworthy subtype annotation or careful
external validation of clusters. Social synchrony needs stable partner identity
and role. Facial expression, semantic object interaction, EEG coupling, and
finger-level grasping require modalities beyond the released body coordinates.
Longitudinal treatment claims require intervention and follow-up information,
not just repeated recordings.

\paragraph{Reading order and durable evidence.}
Read the original methods with the local schema audit first, then Lemler for
evaluation design and MBaye for interpretable recurrence. Read Koehler and
Manelis-Baram to examine context and measurement validity. Use Move4AS, grasping,
multimodal fusion, and VLM work when choosing whether a later research question
justifies additional data. Citation counts were not reliably refreshed and are
not used to rank evidence. Publisher/author URLs, affiliations, dates, access
failures, and paper hashes are recorded in the Markdown report and manifests.
